% Transcription — fonds Alexandre Grothendieck, shelfmark 98, batch 1 (pages 1-20).
% Established from the facsimile archives/batches/98/batch-01.pdf (a mirror of
% https://grothendieck.umontpellier.fr/98.pdf), per the skill
% transcribe-grothendieck, read from pdftoppm renders at 200-600 dpi (the
% tiles were still being cut). The batch file carries no cover sheet: PDF
% page k is folder page k, and the archivists' pencil numbers bottom-left
% agree wherever legible (« 8 », « 10 », « 11 », « 12 », « 13 », « 15 »,
% « 17 », « 18 », « 19 », « 20 »).
%
% Pass: Opus 5.5 (claude-opus-5-5), 2026-09-24 — first pass, unchecked against the pages by a human.
%
% The folder sits in the inventory group "69-89" « Géométrie et topologie
% combinatoire », whose subject does not match this folder's (as for folder
% 100); the folder has its own date, which \dating{} carries verbatim.
% Nothing is concluded here from the group.
%
% Pages transcribed: 2, 4, 6, 8, 10, 11, 12, 13, 15, 17, 19 (his manuscript)
% and 18 (summarised, see below). Absent:
%   - page 1, a buff card cover blank but for his pencil « Formalisation des
%     Normes » and « Cf aussi SGA 4 XVIII »; the words head the \section and
%     are quoted in a \note (hand.md §5);
%   - pages 3, 5, 7, 9, typed pages « - 39 - », « - 35 - », « - 37 - »,
%     « - 36 - » of a French typescript « n° 329 » on general topology
%     (product and quotient spaces), not by him, scanned upside down as the
%     versos of his sheets, with no mark of his hand found and no bearing on
%     norms: skipped under the folder 100/71 rule;
%   - pages 14, 16, 20, typed pages « - 99 - » (n° 266), « - 96 - » and
%     « - 89 - » (n° 268) of a French typescript on abelian varieties
%     (theorem of the square, divisors, linear systems), not by him, with no
%     mark of his (page 20's script letters are inked in by the typist's
%     side, not annotations): skipped for the same reason.
%   - Page 18 is typed page « - 90 - » of the same n° 268 typescript, upside
%     down, carrying pencil annotations (« Prouver ! », « utile au cas non
%     singulier) », a word replacing a struck « théorème », underlinings).
%     Per the rule it is kept: the typed text summarised in one French \note,
%     the pencil marks in \marginal. They are attributed to him only because
%     nothing on the leaf points elsewhere; a reviewer should confirm the hand.
%     Page 11 is a double leaf photographed open (two sheets side by side).
%
% What the batch is. Pages 2-8 (black ink): a formalisation of norms.
% Page 2 starts from the linear algebra — A a ring, M free of rank n,
% L_A(T^n M) = T^n L_A(M) ⊃ TS^n, endomorphisms of ⊗^n M commuting with the
% symmetries acting on Λ^n M — and passes to X^n_S → Symm^n_S(Y). Page 4 is
% a sheet of scattered diagrams for a norm N^α_{F/T/S} : Γ(G/T) → Γ(F/S)
% attached to a finite T → S of degree r and a symmetric map, with the boxed
% transitivity N^γ_{(G⊗F)/U/S} = N^α_{F/T/S} ∘ N^β_{G/U/T}. Page 6 lists the
% properties wanted of the canonical section Can^φ_F: S → Y — cases r = 0, 1,
% functoriality, additivity, transitivity — and page 8 states the
% transitivity for Z-morphisms, boxed: N^φ_{G/X/Z}(cα') =
% N^χ_{F/Y/Z}(N^ψ_{G/X/Y}(α')). Pages 10-13 (blue ink) axiomatise it:
% « Théorie axiomatique des morphismes … »: a category with finite sums, a
% class N of morphisms stable by identity, composition, base change, sums,
% with quotients (X×_Y…×_Y X)/G for G ⊂ S_r; a degree function (i bis)-(v
% bis) with deg of the quotient = card(E^n/G); canonical morphisms
% Can_{X/Y}: Y → Σ^r(X/Y) with (i ter)-(iv ter). Examples: I, G-sets
% (Can from the fibres of cardinal r); II, finite flat morphisms of schemes
% of rank r; III, analytic and topological variants, the latter via X ×_Γ E
% for a group Γ acting on X. Pages 15, 17, 19 turn to the cohomological
% side: C_G the G-sets with E/G finite, hyperfunctors on C_G and their
% equivalence with G-sets, the boxed spectral sequence (A)
% H^*(G,A) ⇐ H^p(T/e, H^q(A)) with H^q(A)(G/F) = H^q(F,A), then its
% generalisation (B) H^*(−/S,Φ) ⇐ H^p(T/S, H^q(Φ)) (« suite spectrale de
% Leray-Čech ») and the case F' ⊂ F ⊂ G, ending « Tout ceci est à
% généraliser aux hyperfaisceaux quelconques ».
%
% Readings to know. The symbol read \mathrm{Can} on page 6 looks like « Cm »
% there and is settled by page 10; the page-13 and page-15 marginal
% « ; dualement » is a doubtful reading; \underline{\mathcal{H}} renders his
% underlined script H. The connective prose of pages 10-19 is fast and much
% of it is left \ill{}; the formulas read.
\documentclass[11pt,a4paper]{article}
% Le préambule commun à toutes les transcriptions du fonds Grothendieck.
%
% Il tient en une idée : **l'incertitude se compose, elle ne se résout pas.**
% Une transcription qui lisse un mot douteux a détruit la seule chose pour
% laquelle on la faisait — un lecteur doit pouvoir savoir, à chaque endroit,
% ce qui était sur le feuillet et ce qui a été ajouté. Les six macros qui
% suivent sont donc l'appareil critique, et non de la décoration.
%
% Les mêmes commandes sont reconnues par `scripts/render.mjs`, qui produit la
% vue de lecture du site. Une macro ajoutée ici doit l'être aussi là-bas,
% faute de quoi le rendu échouera bruyamment — ce qui est voulu : un rendu
% silencieusement faux est pire qu'un rendu absent.

\ProvidesPackage{grothendieck}[2026/08/08 Transcriptions du fonds Grothendieck]

\usepackage{amsmath,amssymb,amsthm}
\usepackage{mathrsfs}
\usepackage{stmaryrd}
% Les diagrammes commutatifs sont partout dans ce fonds : c'est la langue
% même de la géométrie algébrique de Grothendieck, et une transcription qui
% les rendrait en prose ne transcrirait rien.
\usepackage{tikz-cd}
% Français par défaut — c'est la langue des feuillets — mais l'anglais est
% chargé aussi : la traduction et le résumé sont des documents anglais, et la
% typographie française (espaces avant les deux-points) y serait une faute.
% Ces documents basculent par \selectlanguage{english} en tête de corps.
\usepackage[english,french]{babel}
\usepackage{fontspec}
\usepackage{geometry}
\geometry{a4paper,margin=25mm,marginparwidth=28mm}
\usepackage{marginnote}
\usepackage{xcolor}
% ulem, not soul. `soul`'s \st and \ul are built on a character-by-character
% scan that XeLaTeX's font machinery breaks: the document fails with an
% undefined control sequence rather than degrading. ulem does the same two jobs
% and is Unicode-engine safe. `normalem` stops it hijacking \emph, which the
% transcriptions use for Grothendieck's own emphasis.
\usepackage[normalem]{ulem}
\usepackage{hyperref}
\hypersetup{colorlinks=true,linkcolor=black,urlcolor=black,pdfborder={0 0 0}}

% --- Métadonnées du lot -----------------------------------------------------
% Reprises telles quelles par le script de rendu, qui les affiche en tête de
% la vue de lecture. Elles fixent ce que le fichier prétend couvrir : sans
% elles, une transcription flotte sans point d'attache dans un fonds de seize
% mille feuillets.

\newcommand{\folder}[1]{\gdef\grothFolder{#1}}
\newcommand{\batch}[1]{\gdef\grothBatch{#1}}
\newcommand{\leaves}[2]{\gdef\grothFirst{#1}\gdef\grothLast{#2}}
\newcommand{\foldertitle}[1]{\gdef\grothFoldertitle{#1}}
\gdef\grothFolder{?}\gdef\grothBatch{?}\gdef\grothFirst{?}\gdef\grothLast{?}\gdef\grothFoldertitle{}

% --- Le résumé --------------------------------------------------------------
%
% Il ouvre la lecture modernisée, sous le titre « Résumé », et il est composé
% en petit. Ce n'est pas un ornement : c'est le seul passage du document qui
% ne soit pas des mathématiques, et un lecteur doit pouvoir le reconnaître
% avant de l'avoir lu — puis le sauter s'il connaît déjà le sujet, ou s'y
% arrêter s'il ne le connaît pas. Le filet qui le suit marque la reprise du
% corps.

\newenvironment{resume}
  {\small\setlength{\parskip}{0.45em}}
  {\par\medskip\hrule\medskip}

% --- Mots-clés ---------------------------------------------------------------
%
% En anglais, en clôture du résumé de la lecture modernisée : le vocabulaire
% moderne sous lequel chercher ce que les feuillets construisent. Le manifeste
% du site (`npm run manifest`) les extrait pour étiqueter la cote entière —
% cette ligne est la source unique des tags, il n'y a pas de fichier à côté.
\newcommand{\keywords}[1]{\par\smallskip\noindent
  {\footnotesize\textsc{Keywords} --- \textit{#1}}\par}

% --- L'appareil critique ----------------------------------------------------

% Le numéro de feuillet, en marge. C'est l'ancre qui permet de régler un
% désaccord sur une formule en regardant *un* feuillet plutôt qu'une chemise
% de six cents. Le site s'en sert aussi pour tourner les pages du fac-similé
% quand on fait défiler la transcription.
\newcommand{\leaf}[1]{%
  \marginnote{\footnotesize\color{gray!70}\textsf{#1}}%
  \label{leaf:#1}\ignorespaces}

% Illisible : on ne devine pas. Un mot inventé qui se lit comme les autres est
% le pire résultat possible.
\newcommand{\ill}{\textcolor{red!60!black}{[\,\ldots\,]}}

% Lecture douteuse : le mot est donné, et signalé comme incertain.
\newcommand{\uncertain}[1]{\uline{#1}}

% Ajout de l'éditeur — une lettre manquante, un mot sous-entendu.
\newcommand{\add}[1]{\textcolor{blue!50!black}{[#1]}}

% Biffé par Grothendieck lui-même. Ce qu'il a rayé fait partie du document :
% une rature dit souvent où la pensée a changé de direction.
\newcommand{\struck}[1]{\sout{#1}}

% Note du transcripteur, sur l'état matériel du feuillet.
\newcommand{\note}[1]{\par\smallskip{\small\color{gray!60!black}\textit{[#1]}}\par\smallskip}

% Note marginale de Grothendieck, distincte du corps du texte.
\newcommand{\marginal}[1]{\par\smallskip\noindent\hspace{1em}%
  {\small\color{gray!60!black}#1}\par\smallskip}

% --- Titre ------------------------------------------------------------------

\AtBeginDocument{%
  \begin{center}
    {\large\bfseries Fonds Alexandre Grothendieck --- cote n\textsuperscript{o}~\grothFolder}\\[2pt]
    {\itshape\grothFoldertitle}\\[4pt]
    {\small feuillets \grothFirst--\grothLast\ (lot \grothBatch)}\\[2pt]
    {\footnotesize\color{gray!60!black}Transcription établie d'après le fac-similé numérisé
     par l'Université de Montpellier.\\
     Les lectures incertaines sont soulignées, les illisibles notées [\,\ldots\,],
     les ajouts entre crochets.}
  \end{center}
  \vspace{1em}}

\endinput


\folder{98}
\batch{1}
\pages{1}{20}
\dating{[à partir de 1966-1967]}
\watermark{Édition de démonstration}
\foldertitle{Schémas de Hilbert. Normes : notes manuscrites (s.d.).}

\begin{document}

\section{Formalisation des normes}
\note{titre et renvoi écrits au crayon de sa main sur la couverture cartonnée
(page 1), seuls mots de la feuille : « Formalisation des Normes », et au-dessous
« Cf aussi SGA 4 XVIII ».}

\page{2}
\marginal{Normes}

$A$ un anneau \ill{}

$M$ $A$-module libre de rang $n$

$1_M \in L_A(M)$ \quad d'où \quad $1_A \otimes \cdots$

\struck{$M \otimes = 0$}

$L_A\bigl(\overset{n}{\otimes}_A M\bigr) = \overset{n}{\otimes}_A L_A(M) \supset$ \ill{}

\[
  L_A\bigl(T^n_A(M)\bigr) \simeq T^n_A\bigl(L_A(M)\bigr) \supset TS^n_A\bigl(L_A(M)\bigr)
\]

Endomorphismes de $\overset{n}{\otimes}_A M$ qui commutent aux opérations de
symétrie / Ils opèrent en particulier sur la partie
\struck{$\otimes$} $\overset{n}{\Lambda}_A M$ de $\overset{n}{\otimes}_A M$
stable \uncertain{par} symétrie, d'où —

$X \;\; \overset{n}{\boxtimes}_S X \struck{\ill{}} \longrightarrow Y$,
$\mathrm{Cart}^n_S(X)$, $X \to Y$ \note{flèche verticale, $X$ au-dessus de $Y$.}

$S$, \quad $\overset{n}{\otimes} F$, \quad $B$, \quad $A$
\note{lettres isolées, écrites l'une sous l'autre, sans flèche.}

$X \longrightarrow \overset{n}{\boxtimes} X$ \qquad $\overset{n}{\Lambda}_S F$

\begin{tikzcd}
X & X^n_S \arrow[r] & Z \\
Y \arrow[u] & Y^n_S \arrow[u] \arrow[r] & \mathrm{Symm}^n_S(Y) \arrow[u] \\
S & &
\end{tikzcd}

\note{$S$ est écrit sous $Y$, séparé de lui par un point, sans flèche.}

\page{4}
\note{page de diagrammes épars, sans prose ; ils sont donnés de haut en bas
et de gauche à droite.}

\begin{tikzcd}
G & \\
S \arrow[u, no head] & T \arrow[l]
\end{tikzcd}

\note{une flèche monte de $T$, sans but écrit ; à droite de $T$, une flèche
double vers \struck{$S$}, biffée.}

\[
  N^F_{T/S} : \Gamma(G_T/T) \longrightarrow \Gamma(G/S)
\]

\begin{tikzcd}
\struck{G} & G' \arrow[l] & G \arrow[l, "r"'] \\
& S \arrow[u, no head] & T \arrow[l, no head] \arrow[u, no head]
\end{tikzcd}

\note{à droite de $T$, isolée, la lettre $U$.}

\begin{tikzcd}
G & G^r_S \arrow[l] & G^r_T \arrow[l] \\
S \arrow[u, no head] & & T
\end{tikzcd}

$G^r$ \qquad \struck{$N^{r,}_{T/S}($}

\begin{tikzcd}
F & G \arrow[l, "\alpha(r)"'] \\
S \arrow[u] & T \arrow[l] \arrow[u]
\end{tikzcd}

\[
  N^{\alpha}_{F/T/S}(\varphi)
\]
\[
  N^{\alpha}_{F/T/S} : \Gamma(G/T) \to \Gamma(F\ill{}
\]
\note{la ligne s'arrête au bord droit de la feuille.}

\begin{tikzcd}
X \arrow[d, no head] & Y \arrow[l, "\alpha\ (r)"'] \arrow[d] & Z \arrow[l, "\beta\ (s)"'] \arrow[d] \arrow[ll, bend right=30, "(rs)"'] \\
S & T \arrow[l] & U \arrow[l] \\
& F & G \\
& & G \otimes_T F
\end{tikzcd}

\note{sur les flèches horizontales du haut, $(r)$ et $(s)$ sont écrits
au-dessus, $\alpha$ et $\beta$ au-dessous ; l'arc de $Z$ vers $X$ porte
$(rs)$. Les lettres $F$, $G$ et $G \otimes_T F$ sont écrites sous $T$ et
$U$, sans flèche.}

\[
  \boxed{\;N^{\gamma}_{(G \otimes_T F)/U/S}(\varphi)
  = N^{\alpha}_{F/T/S}\bigl(N^{\beta}_{G/U/T}(\varphi)\bigr)\;}
\]
\note{dans l'indice du membre de gauche, une lettre surchargée après
$(G \otimes_T F)/$.}

\begin{tikzcd}
& X_2 & \\
X & \times \arrow[l] & \\
& X_1 \arrow[d] & Y \arrow[l, "r_1"] \arrow[uul, "r_2"'] \arrow[d] \\
& S & T \arrow[l]
\end{tikzcd}

\note{sous $T$ : « $r_1, r_2$ ». Le signe $\times$ est écrit entre $X_2$ et
$X_1$, une flèche allant de lui vers $X$ : on lit $X \leftarrow X_2 \times
X_1$.}

\page{6}
$\overset{n}{\Lambda} M$

\begin{enumerate}
  \item[(i)] $X^r_S \to Y$ ; \quad $X \to X'$, \quad $Y \to Y'$ ;
  \item[(ii)] $F$ au-dessus de $X \to S$, \quad $X^r_S \xrightarrow{\varphi} Y$ ; $F'$ au-dessus de $X' \to S'$, \quad $X'^r_{S'} \xrightarrow{\varphi'} Y'$.
\end{enumerate}
\note{(i) et (ii) sont réunis par une accolade, dont un trait descend vers
la ligne b) ; les flèches $X \to S$ et $X' \to S'$ sont dessinées
verticalement, $F$ et $F'$ écrits au-dessus de $X$ et $X'$.}

a) Cas où $r = 0$, $r = 1$

b) Fonctorialité en $X, Y, S$

\begin{tikzcd}
X^n_S \arrow[r, "\varphi"] \arrow[d] & Y \arrow[d] & S \arrow[l, "\mathrm{Can}^{\varphi}_F"'] \\
X'^n_{S'} \arrow[r, "\varphi'"'] & Y' & S' \arrow[l, "\mathrm{Can}^{\varphi'}_{F'}"]
\end{tikzcd}

\note{dans le carré de gauche : « comm. par hyp. » ; entre les deux flèches
de droite, un signe d'égalité surchargé et « $\mathrm{Can}^{\varphi', n}$ ».
Le symbole transcrit $\mathrm{Can}$ revient tout au long de la page
(une section $S \to Y$ attachée à $\varphi$ et à $F$) ; il est écrit ici
serré, presque « Cm », et se lit sur la page 10, plus nette
($\mathrm{Can}_u = \mathrm{Can}_{X/Y}$). Au bas du
diagramme, le $S$ est surchargé en $S'$.}

\begin{tikzcd}
Y & Y' \arrow[l] \\
S \arrow[u, "\mathrm{Can}^{\varphi}_F"] & S' \arrow[l] \arrow[u, "\mathrm{Can}^{\varphi'}_{F'}"']
\end{tikzcd}

\note{dans le carré : « comm. », avec un signe $\equiv$.}

c) Additivité en $M$

\begin{tikzcd}
X^{r_1+r_2}_S \arrow[r, "\varphi"] & Y & S \arrow[l, "\mathrm{Can}^{\varphi}_{F_1 \times F_2}"'] \\
X^{r_1}_S \times_S X^{r_2}_S \arrow[u, "\wr\ \mathrm{canon.}"] \arrow[r, "\varphi_1 \times_S \varphi_2"'] & Y_1 \times_S Y_2 \arrow[u, "u"'] & S \arrow[l, "\mathrm{Can}^{\varphi_1}_{F_1} \times \mathrm{Can}^{\varphi_2}_{F_2}"] \arrow[u, no head]
\end{tikzcd}

\note{dans le carré de gauche, « comm. » suivi de mots surchargés et
biffés, \ill{} ; les deux $S$ de droite sont reliés par un double trait
vertical, l'égalité.}

d) Transitivité.

\begin{tikzcd}
G & X \arrow[d, "r"] & G \otimes_Y F \\
F & Y \arrow[d, "s"] & \\
& Z &
\end{tikzcd}

\note{$G$ est écrit contre la flèche $r$, $F$ contre la flèche $s$ ; un trait
oblique porte $rs$ de $X$ à $Z$, sous $G \otimes_Y F$. À droite, derrière
un trait vertical : « Cas où \ill{} ».}

\begin{itemize}
  \item $X^r_Y \xrightarrow{\ \varphi\ } \mathfrak{X}_r$ ($Y$-\uncertain{morphisme})
  \item $Y^s_Z \xrightarrow{\ \psi\ } \mathfrak{Y}_s$ ($Z$-\uncertain{morphisme})
  \item $X^r_Z \xrightarrow{\ \chi\ } \mathfrak{Z}_r$ ($Z$-\uncertain{morphisme})
  \item $\mathfrak{X}^s_Z \xrightarrow{\ u\ } \mathfrak{Z}$ ($Z$-\uncertain{morphisme})
\end{itemize}
\note{les quatre lignes sont encadrées ; deux traits obliques barrent les
parenthèses, et le cadre se ferme sur « \ill{} comm. ». Les lettres
$\mathfrak{X}$, $\mathfrak{Y}$, $\mathfrak{Z}$ transcrivent des capitales
cursives ; l'exposant de la troisième ligne est lu $r$, peut-être $rs$.}

\begin{tikzcd}
X^{rs}_Z \arrow[r, "\chi"] & \mathfrak{Z} \\
(X^r_Y)^s_Z \arrow[u, "\mathrm{can}"] \arrow[r, "\varphi^s_Z"'] & \mathfrak{X}^s_Z \arrow[u, "u"']
\end{tikzcd}

\note{ce carré est entouré d'un trait qui le relie à « \uncertain{Fixés}
\ill{} », en haut à gauche, et à « Commutativité », en bas.}

\page{8}
$G^{rs}_Z \xrightarrow[\text{symm}]{\ \varphi\ } G'$ \quad d'où diagramme comm.

\begin{tikzcd}
(G_Y)^{rs}_Y \arrow[r, "\varphi_Y"] \arrow[d, no head] & G'_Y \arrow[d] \\
G^{rs}_Z \arrow[r, "\varphi"'] & G'
\end{tikzcd}

\note{contre la flèche verticale de droite, un signe \ill{}.}

On \uncertain{suppose} \uncertain{qu'il} \ill{} ainsi

$(G_Y)^s \xrightarrow[\text{symm}]{\ \psi\ } G''$ \quad d'où

\struck{$(G_Y$\ill{}}

$G''^{s}_Z \xrightarrow[\text{symm}]{\ \chi\ } G'$ \qquad $Z$-morphisme
\uncertain{rendant} commutatif

\begin{tikzcd}
(G_Y)^{rs} \arrow[r, "\psi^s_Y"] \arrow[dr, "\mathrm{can}\,\varphi_Y"'] & G''^{s}_Z \arrow[d, "\chi"] \\
& G'
\end{tikzcd}

\note{sous $(G_Y)^{rs}$, une flèche verticale biffée.}

Conclusion. $\alpha : X \to G$ un $Z$-morphisme, d'où
\[
  N^{\varphi}_{G/X/Z}(\alpha) : Z \longrightarrow G'
\]
\marginal{\ill{} $\alpha = c\alpha'$}
Mais aussi un $Y$-morphisme $\alpha' : X \to G_Y$, d'où
\[
  N^{\psi}_{G/X/Y}(\alpha') : Y \longrightarrow G''
\]
et \ill{}, en \uncertain{appliquant} $\chi$, on \uncertain{trouve}
\[
  N^{\chi}_{F/Y/Z}\bigl(N^{\psi}_{G/X/Y}(\alpha')\bigr) : Z \longrightarrow G'
\]
Je dis que
\[
  \boxed{\;N^{\varphi}_{G/X/Z}(c\alpha') = N^{\chi}_{F/Y/Z}\bigl(N^{\psi}_{G/X/Y}(\alpha')\bigr)\;}
\]
\note{dans le cadre, l'indice $F/Y/Z$ est écrit par-dessus un premier
$F/X/Z$.}

\page{10}
\subsection{Théorie axiomatique des morphismes \ill{}}

Soit $\mathcal{C}$ une catégorie, où les sommes finies existent, et
soit $\mathcal{N}$ un ensemble de \struck{\ill{}} morphismes de
$\mathcal{C}$, ayant les propriétés suivantes
\note{« où les sommes finies existent » et « ayant les propriétés
suivantes » sont ajoutés au-dessus de la ligne.}
\begin{enumerate}
  \item[(i)] Si $X \in \mathcal{C}$, alors $\mathrm{id}_X \in \mathcal{N}$
  \item[(ii)] Si $u, v \in \mathcal{N}$ sont composables, alors $vu \in \mathcal{N}$
  \item[(iii)] Si $u : X \to Y$ est dans $\mathcal{N}$, et si $Y' \to Y$ est
  quelconque, alors $X' = X \times_Y Y'$ existe et
  $u' : \struck{X \times_Y} X' \to Y'$ est $\in \mathcal{N}$
  \item[(iv)] Si $u_i : X_i \to Y$ ($i = 1, 2$) sont $\in \mathcal{N}$, alors
  $(u_1, u_2) : X_1 \amalg X_2 \to Y$ est dans $\in \mathcal{N}$ ; de plus,
  l'extension de \ill{}
  \note{fin de (iv) très corrigée : mots biffés, ajouts interlinéaires,
  \ill{}.}
  \item[(v)] Si $u : X \to Y$ est $\in \mathcal{N}$, alors \ill{} tout
  sous-groupe $G$ de $\mathfrak{S}_r$, le quotient
  $\bigl(\overbrace{X \times_Y \cdots \times_Y X}^{r}\bigr)/G$ existe,
  \uncertain{et il est} \ill{}
  \note{ces derniers mots sont entourés au crayon rouge, avec un « ? ».}
\end{enumerate}

Soit $\deg$ une fonction sur $\mathcal{N}$, \ill{} : $\deg(X/Y)$
ayant les propriétés suivantes
\begin{enumerate}
  \item[(i bis)] $\deg(\mathrm{id}_X) = 1$
  \item[(ii bis)] $\deg(vu) = \deg v \deg u$
  \item[(iii bis)] $\deg(u) = \deg(u^{Y'})$
  \item[(iv bis)] $\deg(u, v) = \deg u + \deg v$
  \item[(v bis)] Le quotient envisagé par (v) a \ill{} degré qui \uncertain{se}
  calcule ainsi : Soit \struck{$\deg(X/Y)$, \ill{}}
  $r = \deg(X/Y)$, \ill{} $E = [1, r]$,
  \struck{$E = \Lambda^r$, $F = (\overset{r}{\otimes} E)/G$}
  Soit $F = E^n/G$, alors \struck{$\deg((X/Y)^n/G)$} le degré de
  $\bigl(\underbrace{X \times_Y \cdots \times_Y X}_{n}\bigr)/G$ est égal à
  $\mathrm{card}(E^n/G)$.
\end{enumerate}
\note{à la place de « Soit », un premier essai encadré et biffé, \ill{}.}

Enfin, pour \struck{\ill{}} morphisme $u \in \mathcal{N} : X \to Y$
de degré $r$, on se donne un morphisme
\[
  \mathrm{Can}_u = \mathrm{Can}_{X/Y} : Y \to \Sigma^r X
\]
(\struck{\ill{}} \uncertain{canonique} de $u$, ou de $X/Y$), avec les
propriétés suivantes
\note{« de degré $r$ » est ajouté au-dessus de la ligne, relié par une
accolade ; un trait de crayon rouge part de là vers la marge.}
\begin{enumerate}
  \item[(i ter)] $\mathrm{Can}_{X/X} = \mathrm{id}_X$
  \item[(ii ter)] \ill{} commutativité dans le diagramme suivant :
\end{enumerate}

\page{11}
\note{deux feuillets photographiés côte à côte ; celui de gauche continue
directement la page 10, celui de droite ouvre l'Exemple I, que la page 12
poursuit.}

(où $X \xrightarrow[r]{u} Y \xrightarrow[s]{v} Z$ sont des morphismes
$\in \mathcal{N}$, de degrés $r, s$) :
\note{à droite de cette ligne, isolé : $Y \longrightarrow \Sigma^r(X/Y)$.}

\begin{tikzcd}[column sep=small, nodes={font=\scriptsize}]
Z \arrow[r, "\mathrm{Can}_{Y/Z}"] \arrow[drr] \arrow[ddrr, "\mathrm{Can}_{X/Z}"'] & \Sigma^s(Y/Z) \arrow[r, "\Sigma^s(\mathrm{Can}_{X/Y}/Z)"] & \Sigma^s(\Sigma^r(X/Y)/Z) \arrow[d, "\mathrm{can}"] \\
& & \Sigma^s(\Sigma^r(X/Z)/Z) \arrow[d, "\mathrm{can}"] \\
& & \Sigma^{sr}(X/Z)
\end{tikzcd}

\note{sous l'étiquette $\mathrm{Can}_{Y/Z}$, un premier $\mathrm{Can}_{X/Z}$
biffé ; le terme du milieu de la colonne de droite est écrit par-dessus un
premier terme, peut-être $\Sigma^{s+r}$. Les deux flèches obliques issues de $Z$ sont
chacune doublées d'un trait ondulé.}

(iii ter) Soit $u : X \to Y$ dans $\mathcal{N}$, et $Y' \to Y$
  quelconque, alors $Y' \xrightarrow{\mathrm{Can}_{X'/Y'}}$ est « l'image
  inverse » de $\mathrm{Can}_{X/Y}$ par $Y' \to Y$, i.e. le diagramme

\begin{tikzcd}
\Sigma^r(X/Y) & Y \arrow[l, "\mathrm{Can}_{X/Y}"'] \\
\Sigma^r(X'/Y') \arrow[u, "\mathrm{can}"] & Y' \arrow[l, "\mathrm{Can}_{X'/Y'}"] \arrow[u]
\end{tikzcd}

est commutatif

(iv ter) Soient $u_i : X_i \to Y$ ($i = 1, 2$) dans $\mathcal{N}$.
  Soit $X = X_1 \amalg X_2$, $u = (u_1, u_2)$ ; donc pour tout
  \struck{entier} $r$, on a
  \[
    \Sigma^r(X/Y) = \struck{\ill{}} \coprod_{\substack{i+j=r\\ i,j \geqslant 0}}
    \Sigma^i(X_1/Y) \times \Sigma^j(X_2/Y)
  \]
  en particulier, si $r_i = \deg u_i$ et $r = r_1 + r_2 = \deg u$, on a un
  morphisme d'\ill{} canonique
  $\Sigma^{r_1}(X_1/Y) \times \Sigma^{r_2}(X_2/Y) \to \Sigma^r(X/Y)$.
  Ceci dit, on a commutativité dans

\begin{tikzcd}
Y \arrow[r, "{(\mathrm{Can}_{X_1/Y}, \mathrm{Can}_{X_2/Y})}"] \arrow[dr, "\mathrm{Can}_{X/Y}"'] & \Sigma^{r_1}(X_1/Y) \times \Sigma^{r_2}(X_2/Y) \arrow[d] \\
& \Sigma^r(X/Y)
\end{tikzcd}

\note{« d'\ill{} » est ajouté au-dessus de la ligne.}

\textbf{Exemple I}. Soit $G$ un groupe, et $\mathcal{C}$ la catégorie des
ensembles où $G$ opère. Un morphisme $u : X \to Y$ en $\mathcal{C}$ sera
\struck{\ill{} dans $\mathcal{N}$ s'il existe \ill{}} dit de degré $r$ si
pour tout $y \in Y$, \struck{\ill{}} on a $\mathrm{card}(u^{-1}(y)) = r$.
Un morphisme est dans $\mathcal{N}$ s'il existe $r$ tel qu'il soit de degré
$r$. Les conditions (i) à (v) et (i bis) à (v bis) sont trivialement
vérifiées. [N.B. \struck{\ill{}} de produits fibrés et \ill{} : un \ill{}
quelconques existent dans $\mathcal{C}$, \ill{} de même le quotient d'un
objet de $\mathcal{C}$ par un groupe d'automorphismes $G$]. Soit
$u : X \to Y$ de degré $r$. Montrons que $\Sigma^r(X/Y)$ est
\ill{} l'ensemble \struck{\ill{}} réunion des $\Sigma^r(u^{-1}(y))$
($y \in Y$) ; comme $u^{-1}(y)$ a exactement $r$ éléments,
\struck{\ill{}} il y a dans $\Sigma^r(u^{-1}(y))$ un élément privilégié,
classe de la \ill{} de ses éléments rangés dans un ordre quelconque). D'où
une application canonique
\[
  Y \to \Sigma^r(X/Y),
\]
qui évidemment \ill{} aux morphismes de 2 situations, donc \ill{}
\ill{} sur un \struck{\ill{}} \ill{} de $\mathcal{C}$. Si \ill{}
\note{les deux dernières lignes du feuillet droit, au bord de la feuille,
sont en partie hors de l'image.}

\page{12}
le diagramme pour $\mathrm{Can}_{X/Y}$, on vérifie aussitôt les conditions
(i ter) à (iv ter).

\marginal{\uncertain{Expliquer} \ill{} relation \ill{} $\mathcal{C}$ \ill{}
en général, \ill{} \ill{} \uncertain{géométriques} \uncertain{spéciaux}
\ill{} \ill{}}
\textbf{Exemple II}. Soit $\mathcal{C}$ la catégorie des schémas,
(\ill{} \ill{} \ill{}.) Un morphisme $u : X \to Y$ fini et plat est dit
de degré $r$ s'il correspond à un \struck{\ill{}} faisceau
\struck{\ill{}} \ill{} d'algèbres sur $Y$, loc. libre de
\struck{degré} rang $r$ en tout point. Soit $\mathcal{N}$ l'ensemble des
morphismes qui sont finis et plats et \ill{} \ill{} de degré \ill{}. Les
conditions (i) à (v) et (i bis) à (v bis) se vérifient,
\uncertain{comme} \ill{} \struck{\ill{}} très facilement. De plus, on
peut définir de façon naturelle \struck{des applications} un morphisme
$\mathrm{Can}_{X/Y}$, vérifiant les conditions (i) à (iv).
\note{la parenthèse après « schémas » est ajoutée au-dessus de la ligne ;
« $u : X \to Y$ » l'est aussi. La note marginale, à gauche, est reliée à
l'intitulé de l'exemple.}

\textbf{Exemples III} \struck{\ill{}} (Variantes). En géométrie
analytique, on a une variante évidente. En topologie, prenant les
revêtements \uncertain{non} \uncertain{ramifiés} \ill{} \ill{}, on
trouve une variante analogue ; elle est très voisine de l'exemple I,
\struck{\ill{}} \ill{} précisément, soit $X$ un espace topologique
\uncertain{fixé}, \struck{\ill{}} \ill{} \ill{}

\page{13}
Soit $X \in \mathcal{C}$, avec un groupe $\Gamma$ d'automorphismes
\uncertain{discret} ; \ill{} \ill{} \ill{} pour tout sous-groupe $H$
de $G$, \struck{\ill{} $H$ \ill{}} $X/H$ existe, et que pour
$H' \subset H$ \uncertain{invariant} d'indice fini dans $H$,
\struck{$X/H' \to X/H$ soit de \ill{} \ill{} degré $r$} (\ill{} \ill{}
\ill{} \ill{} \ill{} \ill{}, \ill{} \ill{}) \ill{} \ill{}
$X/H'$ sur \ill{} \ill{} fini $H/H'$ soit \emph{non ramifié}, et
$X/H' \to X/H = (X/H')/(H/H')$ est $\in \mathcal{N}$. Soit
$\mathcal{C}'$ la catégorie des ensembles $E$ où $\Gamma$ opère, tels que
$E/\Gamma$ soit fini. Alors pour \struck{\ill{}} $E \in \mathcal{C}'$,
$X \times_{(\Gamma)} E$ existe, soit $\Phi(E) = X \times_{(\Gamma)} E$.
\note{la lettre du groupe est d'abord $\Gamma$, puis $G$ et $H$ ; on
transcrit comme écrit. Le passage biffé entre parenthèses est entouré d'un
cadre relié par un trait à la ligne précédente.}

$\Phi$ est donc un foncteur $\mathcal{C}' \to \mathcal{C}$,
\ill{} aux sommes finies, et on suppose que $\Phi$ transforme
morphismes de degré $r$ en morphismes de degré $r$ (i.e. fonctions
\uncertain{sur} $\mathcal{N}$). De plus, il transforme
\struck{morphismes} \uncertain{canoniques} en \ill{} \ill{}
\marginal{; \uncertain{dualement} \ill{}}
\note{ce dernier alinéa est marqué d'un trait vertical dans la marge
gauche, auquel la note marginale est reliée. La page s'arrête au bas de la
feuille.}

\page{15}
Soit $G$ un groupe, $\mathcal{C}_G$ la catégorie des ensembles $E$ où $G$
opère \struck{à droite} \uncertain{à gauche}, et tels que $E/G$ soit fini.
\ill{} \ill{} par la catégorie des hyper-foncteurs \ill{} ensembles sur
$\mathcal{C}_G$ (\struck{\ill{}} les épimorphismes dans $\mathcal{C}_G$
\ill{} \ill{} morphismes de descente) est équivalente
\marginal{; \uncertain{dualement}} à la catégorie des ensembles
\struck{\ill{}} où $G$ opère \struck{à droite} \uncertain{à gauche} ; si
$A$ est un tel, $\Phi_A$ l'hyperfoncteur associé, on a
\[
  \Phi_A(X) = \struck{\ill{}} \operatorname{Hom}_G(X, E)
\]
\note{les deux dernières lignes (« la catégorie des ensembles … ») sont
tenues par une accolade dans la marge, où s'écrit la note marginale. Dans
la formule, un premier membre de droite, où se lisent $A$, $X$, $(G)$, est
biffé ; le $X$ du membre de gauche est surchargé. Le membre de droite
porte $E$ et non $A$.}

\struck{Ainsi \ill{}, \ill{} :}
Et si $\Phi$ est un hyperfoncteur, alors l'ensemble à opérateurs associé
est
\[
  A = \Phi(G_s)
\]
où $G_s$ désigne $G$ considéré comme ens. à groupe d'opérateurs :
\ill{} $G$ sur lui-même \ill{} \ill{}
\note{l'indice de $G_s$ est surchargé, peut-être $s$ sur $d$.}

\struck{Soit $F$ un sous-groupe de $G$, et \ill{} $\mathcal{C}_{G;F}$ la
sous-catégorie de $\mathcal{C}_G$ formée des $X \in \mathcal{C}_G$ tels que
$F$ \uncertain{opère} \uncertain{trivialement} ; c'est une sous-catégorie
stable par les opérations habituelles [sommes, \ill{}, produits fibrés,
quotients et \ill{} \ill{}].}
\note{ce dernier alinéa est barré de trois grands traits obliques.}

\page{17}
\[
  \text{(A)}\qquad
  \boxed{\;H^*(G, A) \Longleftarrow H^p\bigl(T/e, \underline{\mathcal{H}}^q(A)\bigr)\;}
  \qquad (T \neq \emptyset)
\]
\marginal{\uncertain{s'obtient} \ill{} la suite spectrale de Leray d'un
\emph{recouvrement} d'un \uncertain{objet}}
\note{la note marginale est écrite en oblique dans la marge gauche, contre
le cadre ; le « A » est cerclé.}

où $\underline{\mathcal{H}}^q(A)$ est l'hyperfoncteur sur $\mathcal{C}_G$
\struck{défini par $\underline{\mathcal{H}}^q(A)(G/F) = H^q(F, A)$}
dont la valeur sur $S \in \mathcal{C}_G$ est donnée par
\[
  \underline{\mathcal{H}}^q(A)(S) = R^q_A\bigl(\struck{A}\,
  \struck{\operatorname{Hom}_{(G)}}(S, A)\bigr)
\]
[et \ill{} donc défini \ill{} par les conditions de \struck{\ill{}}
\ill{} \ill{} \uncertain{finies} \ill{} précédente, et \ill{} \ill{}
pour $S = F\backslash G$, on a
\[
  \underline{\mathcal{H}}^q(A)(F\backslash G) = H^q(F, A)\;]
\]
\note{dans la formule, $\operatorname{Hom}$ et son indice $(G)$ sont
barrés en même temps qu'un premier $A$ ; ce qui reste lisible est
$R^q_A(\ldots(S, A))$.}

\ill{} : \uncertain{utilisé} \ill{} les faits suivants
\begin{enumerate}
  \item[(i)] $T \to e$ est un morphisme de descente, \ill{} pour un
  hyperfoncteur, \ill{} \ill{} \struck{\ill{}} \ill{}
  $\Phi$, on a $\Phi(e) = H^0(T/e, \Phi)$
  \item[(ii)] Si $\Phi$ est \uncertain{un hyperfoncteur} injectif, alors
  $H^p(T/e, \Phi) = 0$ pour $p > 0$ (il suffit de prouver
  $\Phi = \Phi_A$, où $A = \operatorname{Hom}_{\mathbb{Z}}(\mathbb{Z}(G), M) = M^G$,
  alors $\Phi(S) = $ \ill{} \uncertain{applications} de $S$ \ill{} $M$ ;
  donc
  $H^p(T/e, \Phi) = H^p($complexe de \ill{} \ill{} défini par $T$ ;
  $M) = 0$ pour \ill{} \ill{} $p > 0$.
\end{enumerate}
\note{« un hyperfoncteur » est ajouté au-dessus de « injectif », relié par
une accolade. La parenthèse ouverte en (ii) ne se ferme pas.}

\page{18}
\note{feuillet d'un tapuscrit qui n'est pas de lui, retourné à l'envers
au dos de la page 17 : page « - 90 - » d'un texte numéroté « n° 268 » sur
les variétés abéliennes. Il achève la preuve d'un théorème (une base
$f_0, \ldots, f_n$ de $L(D)$, développements en série de puissances au
point $P$, rang de la matrice des coefficients, d'où un diviseur $X$
rationnel sur $K$, « Ceci démontre notre théorème »), puis énonce un
théorème : soit $D$ un diviseur positif sur un produit $A \times V$ d'une
variété de groupe commutatif et d'une variété $V$ complète sans
singularités de codimension 1, rationnel sur un corps de définition $k$
tel que $V$ possède un point simple $P$ rationnel sur $k$ ; alors il
existe une variété de groupe commutatif $B$ et un … — l'énoncé
continue au feuillet suivant, absent du lot. Seules les marques au crayon
sont transcrites ; elles sont d'une autre encre que la page 17, et rien
sur la feuille ne les attribue à une autre main.}

\marginal{\ill{}} \note{écrit à la place du mot « théorème », biffé au
crayon dans « Ceci démontre notre théorème ».}

\marginal{utile au cas non singulier)}
\note{en regard des lignes « tés de codimension 1 » et « lequel $D$ soit
rationnel », dont les mots « complète sans singularités de codimension 1 »
et « point simple $P$ rationnel » sont soulignés au crayon.}

\marginal{Prouver !}
\note{sous l'énoncé, relié par un trait à « nel sur $k$ ».}

\page{19}
Plus généralement \struck{\ill{}} posons
\[
  \underline{\mathcal{H}}^q(\Phi)(S) = H^q(-/S, \Phi) =
  \left\{\begin{array}{l}
    \text{valeurs en } \Phi \text{ des foncteurs dérivés } q\text{-ièmes}\\
    \text{de } \Phi \to \Phi(S)
  \end{array}\right.
\]
\marginal{\uncertain{identités} ; \uncertain{dualement}}
\note{la note marginale est reliée par une flèche à la formule. Sous le
signe $=$, une flèche renvoie à un ajout interlinéaire : « définition
\struck{valeurs en $\Phi|S$} des foncteurs dérivés de $\Gamma S$ sur la
catégorie des \struck{\ill{}} \ill{} \ill{} sur $S$ ; il n'intervient
\uncertain{Čech} \ill{} ».}

On montre que
\begin{enumerate}
  \item[(i)] $\underline{\mathcal{H}}^q(\Phi)$ est un
  hyper-préfaisceau qui \struck{\ill{}} \ill{} \uncertain{finies}
  \ill{} produits
  \item[(ii)] $\underline{\mathcal{H}}^q(\Phi_A)(G/F) = H^q(F, A)$
\end{enumerate}
\note{(i) et (ii) sont tenus par une accolade à gauche.}

Sous ces conditions on a la suite spectrale suivante, si $T \to S$ est
\emph{surjectif} (dans $\mathcal{C}_G$)
\[
  \text{(B)}\qquad
  \boxed{\;H^*(-/S, \Phi) \Longleftarrow H^p\bigl(T/S, \underline{\mathcal{H}}^q(\Phi)\bigr)\;}
\]
\marginal{suite spectrale de Leray-\uncertain{Čech}}
\note{le « B » est cerclé ; la note, à droite, est encadrée.}

Si $F' \subset F \subset G$, on peut prendre
$T = G/F'$, $S = G/F$ et on obtient, si $A$ est un $G$-module,
\[
  H^*(F, A) \Longleftarrow H^p\bigl(F \bmod F' ; \underline{\mathcal{H}}^q(A)\bigr)
\]
\note{les deux $A$ de cette formule sont surchargés ; un exposant illisible
suit le premier $F$ du second membre.}

La suite spectrale est la \ill{} \ill{} \ill{}, \uncertain{utilisée} sur
$A$ la structure de $G$-module, et ne retient que celle de $F$-module
\ill{}. La suite spectrale (B) \ill{} \ill{} \ill{} plus générale que (A).

Tout ceci est à généraliser aux hyperfaisceaux quelconques - - - -
\note{la phrase est marquée de trois traits verticaux dans la marge
gauche ; elle finit sur une ligne de tirets, et la page s'arrête là.}

\end{document}
